\documentclass[main.tex]{subfiles}
\begin{document}

\subsection{Problem 1}
\begin{figure}[h]
    \centering
    \includegraphics[width=0.5\linewidth]{images/2021/Q1-2021.png}
\end{figure}
Write the Chisel code for the circuit in Figure 1. The register shall be 8 bits wide and shall
be initialized to 0 on reset. Explain in your own words what the circuit does.

\subsubsection{Solution}
\begin{adjustwidth}{2em}{0pt}
    \begin{verbatim}
        import chisel3._

        class Circuit extends Module {
            val io = IO(new Bundle {
                val din  = Input(UInt(8.W))
                val load = Input(Bool())
                val add  = Input(Bool())
                val dout = Output(UInt(8.W))
            })
        
            // Registers initialized to 0
            val regA = RegInit(0.U(8.W))
            val regB = RegInit(0.U(8.W))
        
            // regA load behavior
            regA := Mux(io.load, io.din, regA)
        
            // Adder
            val sum = regA + regB
        
            // regB behavior
            regB := Mux(io.add, sum, 0.U)
        
            // Output
            io.dout := regB
        }
    \end{verbatim}
\end{adjustwidth}

\newpage
\subsection{Problem 2}
Given following Chisel code:
\begin{verbatim}
    val reg = RegInit(0.U(8.W))
    switch(sel) {
        is(0.U) { reg := 0.U}
        is(1.U) { reg := reg}
        is(2.U) { reg := reg << din}
        is(3.U) { reg := reg + din}
        is(4.U) { reg := reg - din}
        is(5.U) { reg := reg + reg}
    }
\end{verbatim}
Explain the function and draw a block diagram.

\subsubsection{Solution}
\begin{figure} [h]
    \centering
    \includegraphics[width=0.5\linewidth]{images//2021/Q2-2021.png}
\end{figure}
\begin{adjustwidth}{2em}{0pt}
    This circuit consists of a register connected to a MUX controlled by a select signal Sel. This allows for the operations:
    \begin{itemize}
        \item Set Reg = 0.
        \item Set Reg = Reg.
        \item Set Reg = Reg + din.
        \item Set Reg = Reg - din.
        \item Set Reg = Reg + Reg.
        \item And left-shift Reg by din bits.
    \end{itemize}
\end{adjustwidth}

\newpage
\newpage
\subsection{Problem 3}
To store state in a digital circuit we need a form of memory. In an ASIC and in an FPGA there are two types of memories: flip-flops and static RAM (SRAM) blocks. What is the cost in number of transistors for an SRAM cell and a flip-flop? Find the (English) Wikipedia entries of SRAM and a master-slave D flip-flop and explore the schematics. How
many transistors are needed to implement CMOS SRAM cell or a master-slave D flip-flop? (remember, a NAND or NOR gate needs 4 transistors, an inverter 2 transistors

\subsubsection{Solution}
\begin{adjustwidth}{2em}{0pt}

    \note{This solution was provided by ChatGTP, not at calculations/solutions might be correct.}
    
    
    To compare the implementation cost of memory in digital circuits, we count the number of transistors required for a CMOS SRAM cell and a master-slave D flip-flop.
    
    \subsubsection*{SRAM Cell}
    
    A standard CMOS SRAM bit cell is implemented as a \textbf{6-transistor (6T) cell}. It consists of:
    
    \begin{itemize}
        \item Two cross-coupled inverters used to store the bit
        \item Two access transistors used for read/write access
    \end{itemize}
    
    Since an inverter requires 2 transistors:
    
    \[
    2 \times 2 = 4
    \]
    
    transistors are needed for the two inverters.
    
    Additionally, there are two access transistors:
    
    \[
    4 + 2 = 6
    \]
    
    Thus, a CMOS SRAM cell requires:
    
    \[
    \boxed{6 \text{ transistors per bit}}
    \]
    
    \subsubsection*{Master-Slave D Flip-Flop}
    
    A master-slave D flip-flop consists of two latches:
    
    \[
    \text{Master latch} + \text{Slave latch}
    \]
    
    Each latch can be implemented using:
    
    \begin{itemize}
        \item Two inverters
        \item Two transmission gates
    \end{itemize}
    
    Since:
    
    \[
    1 \text{ inverter} = 2 \text{ transistors}
    \]
    
    and
    
    \[
    1 \text{ transmission gate} = 2 \text{ transistors}
    \]
    
    the transistor count for one latch is:
    
    \[
    (2 \times 2) + (2 \times 2) = 8
    \]
    
    Since a master-slave flip-flop uses two latches:
    
    \[
    2 \times 8 = 16
    \]
    
    Thus, a master-slave D flip-flop requires:
    
    \[
    \boxed{16 \text{ transistors per bit}}
    \]
    
    \subsubsection*{Conclusion}
    
    Comparing the two memory implementations:
    
    \[
    \text{SRAM cell} = 6 \text{ transistors/bit}
    \]
    
    \[
    \text{Master-slave D flip-flop} = 16 \text{ transistors/bit}
    \]
    
    SRAM is therefore significantly more area-efficient, which is why large memories are implemented using SRAM cells, while flip-flops are mainly used for registers and sequential logic.
\end{adjustwidth}

\newpage
\subsection{Problem 4}
Draw the state machine and datapath for the following circuit:
\begin{itemize}
    \item A vending machine that accepts coins of 5 and 10 kr.
    \item A soda costs 10 kr. When the sum of coins inserted into the vending machine is larger or equal to 10 kr.
    \item One can buy a soda by pressing the buy button.
    \item The inputs to the circuit are coin5, coin10, buy, and return. 
    \item By pressing the return button, any money in the machine will be returned to the customer.
    \item Assume that all input signals are valid for one clock cycle (no edged detection needed).
    \item The outputs are a release and a moneyReturn signal.
    \item The output signals have to be valid for only one clock cycle.
\end{itemize}
Note, that this exercise is related to your vending machine project, but it is NOT your vending machine. The specification is different (considerably simpler)

\subsubsection{Solution}
\begin{figure} [h]
    \centering
    \includegraphics[width=0.5\linewidth]{images//2021/Q4-state-2021.png}
\end{figure}
\begin{figure} [h]
    \centering
    \includegraphics[width=0.7\linewidth]{images/2021/Q4-datapath-2021.png}
\end{figure}

\newpage
\subsection{Problem 5}
Find and describe the critical path in the circuit shown in the figure.
What is the maximum clock frequency?
\begin{multicols}{2}
    \centering
    \includegraphics[width=1\linewidth]{images//2021/Q5-circuit-2021.png}
    \columnbreak
    \includegraphics[width=1\linewidth]{images//2021/Q5-timing-2021.png}
\end{multicols}

\subsubsection{Solution}
\begin{adjustwidth}{2em}{0pt}
    The critical path is the longest register-register path. Looking at the circuit we can find the longest path to be through, R1-G1-G2-G3-G7-R6 and equally as long, R1-G1-G2-G4-G5-R5. 

    Calculating the time we get, DFF clock-to-Q delay + 4 gate delays + setup time. With tCQ prop = 11ps, gate prop = 7ps, tSU = 3ps.
    \begin{align}
        11\text{ ps} + 4 \cdot 7\text{ ps} + 3\text{ ps} = 42\text{ ps}
    \end{align}

    From this we can find the maximum clock frequency with: 
    \begin{align}
        f_{\max} = \frac{1}{T_{\min}} = \frac{1}{42 \cdot 10^{-12}\text{ s}} = 23.8\text{ GHz}
    \end{align}
\end{adjustwidth}

\end{document}